\clearpage
\par\bigskip\pdfbookmark[0]{Über den Zusammenhang. Fragment III}{nexus-fragment-III}\centerline{[III.]}\medskip
Es sei $C$ die Gesamtheit der Repräsentanten sämtlicher Klassen der eigentlich primitiven Formen für die Determinante $-D$. Wir bezeichnen mit $(n)$ die Anzahl aller Darstellungen der Zahl $n$ durch Formen aus der Gesamtheit $C$. Es sei $p$ eine ungerade Primzahl. Dann ist

\SourceLeaf{293}
\[
\begin{aligned}
&1.\quad(pn)=(n)\quad\text{wenn }p\text{ ein Divisor von }D\\
&\left.\begin{aligned}2.\quad(pn)&=(n)+(h)\\3.\quad(pn)&=-(n)+(h)\end{aligned}\right\}
\ \text{wenn }p\text{ Nichtdivisor von }D\\
&\hspace{7em}\left\{\begin{aligned}&\text{Divisor von }xx+D\\&\text{Nichtdivisor von }xx+D\end{aligned}\right.
\end{aligned}
\]
wobei $n=hp^\mu$, $\mu$ beliebig und $h$ nicht durch $p$ teilbar ist.
\[
\begin{aligned}
\text{Im Fall }1.\ \text{ist}\quad&(h)=(ph)=(pph)=(p^3h)\text{ usw.}\\
2.\quad&(ph)=2(h),\ (pph)=3(h),\ (p^3h)=4(h)\text{ usw.}\\
3.\quad&(ph)=0,\ (pph)=(h),\ (p^3h)=0,\ (p^4h)=(h)\text{ usw.}
\end{aligned}
\]
\begin{center}$\ast\qquad\ast$\par$\ast$\end{center}

Aus jeder eigentlich primitiven Klasse für die Determinante $=-D$, deren Anzahl $=\lambda$ ist, sei eine Form ausgewählt; die Gesamtheit dieser Formen sei $L$.

Man bezeichne mit $fA$ die Anzahl sämtlicher Darstellungen der Zahl $A$ durch Formen aus $L$.

Es sei ferner $f(A;p)=f\frac A{p^\alpha}$, wenn $p^\alpha$ die höchste Potenz der Primzahl $p$ ist, die $A$ teilt; ferner $f(A;p,q)=f\frac A{p^\alpha q^\beta}$, wenn $q$ eine andere Primzahl ist, deren höchste $A$ teilende Potenz $=q^\beta$ ist; ebenso $f(A;p,q,r)=f\frac A{p^\alpha q^\beta r^\gamma}$, wenn $r$ eine dritte Primzahl ist, deren höchste $A$ teilende Potenz $r^\gamma$ ist, und so weiter.
\BookRule{.30\textwidth}
